\chapter{数的几何}\label{chap:nt-5}
\section{Euclid格与Minkowski定理}\label{sec:nt-5-1}
\subsection{定义}\label{sec:nt-5-1-1}
一个\newterm{nt-euclidean-lattice}{Euclid格}{Euclidean lattice}由以下资料组成：
\begin{itemize}
\item 一个有限秩自由Abel 群 $L$；
\item $V=\mathbb R\otimes_{\mathbb Z}L$ 上的一个Euclid内积 $\langle\cdot,\cdot\rangle$.
\end{itemize}
设 $L$ 的秩为 $n$.商空间 $V/L\simeq(\mathbb R/\mathbb Z)^n$ 是紧的：此时称 $L$ 是\newterm{nt-cocompact}{余紧的}{cocompact}.
\begin{exercise*}\phantomsection\label{ex:nt-5-1-lattice}
设 $(V,\langle\cdot,\cdot\rangle)$ 为有限维Euclid空间，$L$ 为 $V$ 的一个子群.则下列条件等价：
\begin{enumerate}[label=\roman*)]
\item $L$ 是 $V$ 中的一个（Euclid）格；
\item $L$ 在 $V$ 中离散且余紧；
\item $L$ 包含 $V$ 的一组基，且对每个有界集 $B$，$B\cap L$ 都是有限集.
\end{enumerate}
\end{exercise*}
若 $(e_1,\ldots,e_n)$ 为 $V$ 的一组标准正交基，从而给出同构 $\mathbb R^n\xrightarrow{\sim}V$，就用 $\mathbb R^n$ 上的Lebesgue测度赋予 $V$ 一个测度.这个测度不依赖于标准正交基的选择（正交换基的Jacobi行列式为 $1$）.\footnote{译注：这里测度变换所用的是Jacobi行列式的绝对值，它等于 $1$；正交矩阵的行列式本身也可能为 $-1$.}
由于投影 $V\longrightarrow V/L$ 是局部同胚，这在商空间 $V/L$ 上诱导出一个测度.又因后者为紧空间，其体积 $\operatorname{vol}(V/L)$ 有限；这个体积称为 $L$ 的\newterm{nt-covolume}{余体积}{covolume}，记为 $\operatorname{covol}(L)$.
\subsection{引理}\label{sec:nt-5-1-2}
设 $v_1,\ldots,v_n$ 为 $L$ 的一组基，则
\begin{align*}
\operatorname{covol}(L)^2=\det(\langle v_i,v_j\rangle)_{i,j}.
\end{align*}
\begin{proof}
按商测度的定义，$L$ 的余体积等于 $L$ 在 $V$ 中的一个基本域 $D$ 的体积 $\operatorname{vol}(D)$（即 $D$ 是一个可测集，满足 $V=\bigsqcup_{x\in L}(x+D)$）.平行多面体
\begin{align*}
D=\prod_{i=1}^n[0,1[v_i=\left\{v=\sum_i\lambda_iv_i\in V,\ 0\le\lambda_i<1\right\}
\end{align*}
显然是一个基本域，其体积为
\begin{align*}
\operatorname{vol}(D)=|\det(P)|=\sqrt{\det({}^tPP)},\qquad P=(\langle e_i,v_j\rangle)_{i,j},
\end{align*}
其中 $P$ 是从 $e_i$ 到 $v_j$ 的换基矩阵.而 ${}^tPP=(\langle v_i,v_j\rangle)_{i,j}$.
\end{proof}
\subsection{定理（Minkowski）}\label{sec:nt-5-1-3}
设 $(L,\langle\cdot,\cdot\rangle)$ 为Euclid格，$B\subset V=\mathbb R\otimes L$ 为 $V$ 中有界、凸且对称的子集，并满足下列两个条件之一：
\begin{enumerate}[label=(\alph*)]
\item $\operatorname{vol}(B)>2^n\operatorname{covol}(L)$；
\item $\operatorname{vol}(B)\ge2^n\operatorname{covol}(L)$，且 $B$ 为闭集.
\end{enumerate}
则 $(B\cap L)\setminus\{0\}\ne\varnothing$.
\begin{proof}
\textup{(a)} 考虑投影 $\pi\colon V\twoheadrightarrow V/2L$，注意到 $2^n\operatorname{covol}(L)=\operatorname{covol}(2L)$.因此，假设(a)意味着 $\pi|_B$ 不是单射，故存在 $x\ne y\in B$，使 $x-y\in2L$.于是 $\frac12(x-y)$ 属于 $B$（由凸性及对称性），也属于 $L$，且非零.
\textup{(b)} 由(a)，对每个 $m\in\mathbb N$，集合 $\bigl(L\cap(1+\frac1m)B\bigr)\setminus\{0\}$ 非空.此外，这些集合都是有限集，且构成递减序列，故其交非空.由于 $B$ 是闭集，这个交中的每个元素都属于 $B$.\footnote{译注：这里含 $1/m$ 的序列从 $m\ge1$ 开始；原文写 $m\in\mathbb N$.}
\end{proof}
\section{类数的有限性}\label{sec:nt-5-2}
设 $K$ 为 $\mathbb Q$ 上次数为 $n$ 的数域.固定 $K$ 到 $\mathbb C$ 的嵌入在共轭作用下的各类的一个代表集 $\Sigma$，记 $\Sigma=\Sigma_1\sqcup\Sigma_2$ 为实嵌入与非实嵌入的划分.回顾分解
\begin{align*}
K_{\mathbb R}:=K\otimes_{\mathbb Q}\mathbb R\simeq\prod_{\sigma\in\Sigma}K_\sigma=\mathbb R^{r_1}\times\mathbb C^{r_2},
\end{align*}
其中 $r_1=|\Sigma_1|$ 为实嵌入的个数，$r_2=|\Sigma_2|$ 为非实复嵌入模共轭的个数.在 $K_{\mathbb R}$ 上赋予一个 $\mathbb R$ 线性对合 $x\mapsto\overline x$，它在每个 $\mathbb R$ 因子上为恒等映射，在每个 $\mathbb C$ 因子上为复共轭.这样，对 $x=(x_\sigma)_{\sigma\in\Sigma}$ 及 $y=(y_\sigma)_{\sigma\in\Sigma}$，可以定义如下内积：
\begin{align*}
\langle x,y\rangle:=\operatorname{Tr}_{K_{\mathbb R}/\mathbb R}(x\overline y)=\sum_{\sigma\in\Sigma}\operatorname{Tr}_{K_\sigma/\mathbb R}(x_\sigma\overline{y_\sigma}).
\end{align*}
注意，$(x,y)\in\mathbb C^2\mapsto\operatorname{Tr}_{\mathbb C/\mathbb R}(x\overline y)$ 是 $\mathbb C$ 上通常内积的两倍，因而可以写成
\begin{align*}
\forall x,y\in K\subset K_{\mathbb R},\qquad \langle x,y\rangle=\sum_{\sigma\colon K\hookrightarrow\mathbb C}\sigma(x)\overline{\sigma(y)}.\tag{$*$}\label{eq:nt-5-2-inner}
\end{align*}
这里的和确实遍历所有嵌入，特别地有
\begin{align*}
\forall x\in K\subset K_{\mathbb R},\qquad \|x\|^2=\sum_{\sigma\colon K\hookrightarrow\mathbb C}|\sigma(x)|^2.\tag{$**$}\label{eq:nt-5-2-length}
\end{align*}
\subsection{命题}\label{sec:nt-5-2-1}
沿用以上记号，再令 $D_K:=\operatorname{disc}(\mathcal O_K/\mathbb Z)\in\mathbb Z$.
\begin{enumerate}[label=\roman*)]
\item $\mathcal O_K$ 在 $K_{\mathbb R}$ 中的像是一个余体积为 $\sqrt{|D_K|}$ 的Euclid格.
\item $\mathcal O_K$ 的分式理想 $I$ 在 $K_{\mathbb R}$ 中的像是一个Euclid格，其余体积为
\begin{align*}
\operatorname{covol}(I)=(\mathcal O_K:I)\sqrt{|D_K|}.
\end{align*}
\end{enumerate}
\begin{proof}
\textup{i)} $\mathcal O_K$ 包含 $K$ 的一组 $\mathbb Q$-基，从而包含 $K_{\mathbb R}$ 的一组 $\mathbb R$-基.对 $r>0$，令 $B_r:=\{x\in K_{\mathbb R},\ \|x\|\le r\}$.由前面的\hyperref[ex:nt-5-1-lattice]{习题}，只须证明 $B_r\cap\mathcal O_K$ 是有限集.由上式\eqref{eq:nt-5-2-length}，
\begin{align*}
x\in\mathcal O_K\cap B_r\ \Longrightarrow\ |\sigma(x)|\le r,\quad\forall\sigma\colon K\hookrightarrow\mathbb C.
\end{align*}
对这样的 $x$，特征多项式
\begin{align*}
\chi_{K/\mathbb Q,x}(T)=\prod_{\sigma\colon K\hookrightarrow\mathbb C}(T-\sigma(x))
\end{align*}
的次数为 $[K:\mathbb Q]$，系数为整数，且由只依赖于 $r$ 的常数控制.这样的多项式只有有限多个，故这样的 $x$ 也只有有限多个.
现在设 $\omega_1,\ldots,\omega_n$ 为 $\mathcal O_K$ 的一组 $\mathbb Z$-基.有 $\operatorname{covol}(\mathcal O_K)^2=\det(\langle\omega_i,\omega_j\rangle)_{i,j}$.引入矩阵 $M:=(\sigma_i(\omega_j))_{i,j}$，其中 $\sigma_1,\ldots,\sigma_n$ 为 $K$ 到 $\mathbb C$ 的各个嵌入.由上式\eqref{eq:nt-5-2-inner}，$\langle\omega_i,\omega_j\rangle=\sum_k\sigma_k(\omega_i)\overline{\sigma_k(\omega_j)}$，所以
\begin{align*}
\operatorname{covol}(\mathcal O_K)^2=\det({}^tM\overline M)=|\det(M)|^2.
\end{align*}
另一方面，前面已经见到
\begin{align*}
\operatorname{disc}(\mathcal O_K/\mathbb Z)=D_{\mathcal O_K/\mathbb Z}(\omega_1,\ldots,\omega_n)=\det({}^tMM)=\det(M)^2.
\end{align*}
由此得到余体积公式 $\operatorname{covol}(\mathcal O_K)=\sqrt{|D_K|}$.
\textup{ii)} 若 $I\subset\mathcal O_K$，则因 $I$ 在 $\mathcal O_K$ 中的指数 $(\mathcal O_K:I)$ 有限，ii)由i)推出.\footnote{译注：原文只证明了整理想 $I\subset\mathcal O_K$ 的情形，未写一般分式理想的归约步骤.}
\end{proof}
\begin{lemma*}\phantomsection\label{lem:nt-5-2-bounded-index}
若 $m\in\mathbb N$，则 $\mathcal O_K$ 中指数不超过 $m$ 的理想 $I$ 只有有限多个.
\end{lemma*}
\begin{proof}
若 $(\mathcal O:I)\le m$，则 $m!$ 零化有限群 $\mathcal O_K/I$，所以 $I\supset m!\mathcal O_K$.而 $\mathcal O_K/m!\mathcal O_K$ 是有限集.
\end{proof}
\subsection{定理}\label{sec:nt-5-2-2}
对 $\mathcal O_K$ 的每个分式理想 $J$，都存在理想 $I\subset\mathcal O_K$，使得 $IJ$ 为主理想，且
\begin{align*}
(\mathcal O_K:I)\le\left(\frac2\pi\right)^{r_2}\sqrt{|D_K|}.
\end{align*}
在证明定理之前，先给出我们所关心的推论；它还用到前面的\hyperref[lem:nt-5-2-bounded-index]{引理}.
\begin{corollary*}\phantomsection\label{cor:nt-5-2-class-finite}
类群 $\mathcal C\ell(\mathcal O_K)$ 是有限群.
\end{corollary*}
\begin{proof}[定理的证明]
先重新解释这个陈述.所求理想 $I$ 形如 $(\alpha)J^{-1}$，其中 $\alpha\in K^\times$；由于 $I\subset\mathcal O_K$，必有 $\alpha\in J$.还有
\begin{align*}
(\mathcal O_K:I)&=(\mathcal O_K:\alpha\mathcal O_K)(\mathcal O_K:J)^{-1}\\
&=|N_{K/\mathbb Q}(\alpha)|_\infty\cdot\operatorname{covol}(J)^{-1}\cdot\sqrt{|D_K|}.\tag{$\dagger$}\label{eq:nt-5-2-index}
\end{align*}
因此，问题归结为找出非零的 $\alpha\in J$，使 $|N_{K/\mathbb Q}(\alpha)|_\infty$ 尽可能小.现在注意到
\begin{align*}
|N_{K/\mathbb Q}(\alpha)|_\infty=\prod_{\sigma\in\Sigma}|\sigma(\alpha)|^{[K_\sigma:\mathbb R]}.\tag{$\dagger\dagger$}\label{eq:nt-5-2-norm}
\end{align*}
于是，对 $c=(c_\sigma)_{\sigma\in\Sigma}\in(\mathbb R_+^\times)^\Sigma$，考虑下列球的乘积 $B(c)$：
\begin{align*}
B(c):=\{x=(x_\sigma)_{\sigma\in\Sigma}\in K_{\mathbb R},\ |x_\sigma|\le c_\sigma,\ \forall\sigma\in\Sigma\}.
\end{align*}
这样，$B(c)$ 是闭的、对称的、凸的，其体积为
\begin{align*}
\operatorname{vol}(B(c))&=\prod_{\sigma\in\Sigma_1}(2c_\sigma)\times\prod_{\sigma\in\Sigma_2}(2\pi c_\sigma^2)\\
&=2^n\left(\frac\pi2\right)^{r_2}\prod_{\sigma\in\Sigma}c_\sigma^{[K_\sigma:\mathbb R]},
\end{align*}
其中 $\Sigma=\Sigma_1\sqcup\Sigma_2$ 为实嵌入与非实嵌入的划分.选取 $c$，使得
\begin{align*}
N(c):=\prod_{\sigma\in\Sigma}c_\sigma^{[K_\sigma:\mathbb R]}=\left(\frac2\pi\right)^{r_2}\operatorname{covol}(J). 
\end{align*}
由Minkowski定理~\ref{sec:nt-5-1-3}，$B(c)\cap J$ 中存在非零元素 $\alpha$.对这样的元素，式\eqref{eq:nt-5-2-norm}给出 $|N_{K/\mathbb Q}(\alpha)|\le N(c)$，而式\eqref{eq:nt-5-2-index}给出所需的 $(\mathcal O:I)$ 的上界.
\end{proof}
利用相同方法，并使用算术--几何平均不等式，可以改进定理中 $(\mathcal O:I)$ 的上界所含的常数因子.下列数
\begin{align*}
M_K:=\left(\frac4\pi\right)^{r_2}\frac{n!}{n^n}
\end{align*}
称为\newterm{nt-minkowski-constant}{Minkowski常数}{Minkowski constant}.
\subsection{定理}\label{sec:nt-5-2-3}
对 $\mathcal O_K$ 的每个分式理想 $J$，都存在理想 $I\subset\mathcal O_K$，使得 $IJ$ 为主理想，且 $(\mathcal O_K:I)\le M_K\sqrt{|D_K|}$.
\begin{proof}
算术--几何平均不等式与公式\eqref{eq:nt-5-2-norm}给出
\begin{align*}
|N_{K/\mathbb Q}(\alpha)|_\infty&=\prod_{\sigma\in\Sigma}|\sigma(\alpha)|^{[K_\sigma:\mathbb R]}\\
&\le\left(\frac1n\sum_{\sigma\in\Sigma}[K_\sigma:\mathbb R]|\sigma(\alpha)|\right)^n.
\end{align*}
这提示我们，对 $r>0$，考虑下列有界、对称且闭的集合：
\begin{align*}
B'(r):=\left\{x=(x_\sigma)_{\sigma\in\Sigma},\ \sum_{\sigma\in\Sigma}[K_\sigma:\mathbb R]|x_\sigma|\le r\right\}.
\end{align*}
可以证明，其体积为 $\operatorname{vol}(B'(r))=2^{r_1}\pi^{r_2}\frac{r^n}{n!}$.选取 $r$，使 $\operatorname{vol}(B'(r))=2^n\operatorname{covol}(J)$，于是 $B'(r)\cap J$ 中存在非零元素 $\alpha$.此时
\begin{align*}
|N_{K/\mathbb Q}(\alpha)|_\infty\le\left(\frac rn\right)^n=M_K\cdot\operatorname{covol}(J).
\end{align*}
由前一证明中的等式\eqref{eq:nt-5-2-index}，对理想 $I:=\alpha J^{-1}$，便有 $(\mathcal O_K:I)\le M_K\sqrt{|D_K|}$.
\end{proof}
这个定理对于估计数域的类数当然十分重要；当判别式较小时，甚至可以用它计算类群.不过，先给出一个较为意外的推论：
\begin{corollary*}\phantomsection\label{cor:nt-5-2-ramified-prime}
对每个数域 $K$，至少有一个素数在 $\mathcal O_K$ 中分歧.\footnote{译注：这一断言须排除 $K=\mathbb Q$；原文陈述中未写此限制，而证明最后明确使用了 $n>1$.}
\end{corollary*}
\begin{proof}
只须说明 $D_K\ne\pm1$.出发点是：由于至少存在一个理想类，前一定理保证 $M_K\sqrt{|D_K|}\ge1$.这意味着
\begin{align*}
\sqrt{|D_K|}\ge\left(\frac\pi4\right)^{r_2}\frac{n^n}{n!}\ge\frac{\pi^{n/2}}2>1,
\end{align*}
这里用了 $r_2\le n/2$、$n^n\ge2^{n-1}n!$，最后还用了 $n>1$.
\end{proof}
现在转向更自然的应用.
\begin{example*}\phantomsection\label{ex:nt-5-2-minus-five}
对 $K=\mathbb Q[\sqrt{-5}]$，计算类群 $\mathcal C\ell(\mathcal O_K)$.有 $n=2$、$r_2=1$ 及 $\operatorname{disc}(\mathcal O_K/\mathbb Z)=-20$，所以 $M_K<3$.因此，每个理想类都包含一个在 $\mathcal O_K$ 中指数不超过 $2$ 的代表元 $I$.\footnote{译注：原文这里写 $M_K<3$；推出下一句所需的是 $M_K\sqrt{|D_K|}<3$，本例中此量为 $2\sqrt{20}/\pi<3$.}若 $(\mathcal O_K:I)=1$，则 $I=\mathcal O_K$.
若 $(\mathcal O_K:I)=2$，则 $2\in I$.已知 $2$ 在 $\mathcal O_K$ 中分歧，故存在 $\mathfrak p\in\operatorname{Max}(\mathcal O_K)$，使得 $2\mathcal O_K=\mathfrak p^2$.于是 $(\mathcal O_K:\mathfrak p^2)=4$，从而 $(\mathcal O_K:\mathfrak p)=2$，且 $I=\mathfrak p$.
余下只须确定 $\mathfrak p$ 是否为主理想.然而，$\mathfrak p=(\alpha)$ 蕴含 $|N_{K/\mathbb Q}(\alpha)|=2$.写成 $\alpha=a+b\sqrt{-5}$，则 $|N_{K/\mathbb Q}(\alpha)|=a^2+5b^2$，它不可能等于 $2$.因此，我们已经证明 $\mathcal C\ell(\mathcal O_K)=\mathbb Z/2\mathbb Z$.
\end{example*}
\emph{应用.}说明如何证明Diophantus方程 $y^2+5=x^3$ 没有整数解.设有解 $(x,y)\in\mathbb Z^2$.则在 $\mathcal O_K$ 中有理想分解 $(y+\sqrt{-5})(y-\sqrt{-5})=(x)^3$.由于 $(y+\sqrt{-5})$ 与 $(y-\sqrt{-5})$ 互素，\footnote{事实上，若素理想 $\mathfrak q$ 整除这两个理想，则它整除 $(2\sqrt{-5})$.然而，显然 $(\sqrt{-5})$ 不整除 $(y+\sqrt{-5})$，因为 $y\ne0$；若理想 $\mathfrak p$（满足 $\mathfrak p^2=(2)$）整除 $(y+\sqrt{-5})$，则 $2$ 整除 $x$，从而 $y$ 必为奇数.这是不可能的，因为将 $y^2+5=x^3$ 模 $4$ 约化，就会得到 $1+5=0\pmod4$.}\footnote{译注：原文小字说明以 $y\ne0$ 说明 $(\sqrt{-5})$ 不整除 $(y+\sqrt{-5})$，这一理由不足；这里所需的是排除 $5\mid y$，原文未给出相应论证.}可知存在理想 $I\subset\mathcal O_K$，使得 $(y+\sqrt{-5})=I^3$ 且 $(y-\sqrt{-5})=\overline I^{3}$（在Galois作用下共轭）.由于 $\mathcal C\ell(\mathcal O_K)=\mathbb Z/2\mathbb Z$，$I^3\sim1\Rightarrow I\sim1$，所以 $I$ 是主理想，设 $I=(\alpha)$.于是 $y+\sqrt{-5}=\pm\alpha^3$.但用初等方法即可证明这样的 $\alpha$ 不存在.

\section{Dirichlet 单位定理}\label{sec:nt-5-3}
这里研究数域整数环的单位群 $\mathcal O_K^\times$ 的结构.先给出一个单位判别准则：
\subsection{命题}\label{sec:nt-5-3-1}
设 $\alpha\in\mathcal O_K$.则 $\alpha\in\mathcal O_K^\times$ 当且仅当 $N_{K/\mathbb Q}(\alpha)=\pm1$.
\begin{proof}
若 $\alpha\in\mathcal O_K^\times$，则 $N(\alpha)N(\alpha^{-1})=1$，所以 $N(\alpha)\in\mathbb Z^\times=\{\pm1\}$.反过来，在表达式 $f_\alpha(X)=X^m+a_{m-1}X^{m-1}+\cdots+a_0$ 中，有 $N(\alpha)=\pm a_0$.\footnote{译注：原文此处将 $N_{K/\mathbb Q}(\alpha)$ 写成 $\pm a_0$；一般应为 $((-1)^m a_0)^{[K:\mathbb Q(\alpha)]}$.下面的 $g$ 又将系数下标写成 $a_{n-1}$，与前面次数 $m$ 不一致.这里保留原文两处写法；单位判别结论本身成立.}因此，若 $N(\alpha)=\pm1$，则多项式
\begin{align*}
g(X)=X^m+a_0^{-1}a_1X^{m-1}+\cdots+a_0^{-1}a_{n-1}X+a_0^{-1}
\end{align*}
属于 $\mathbb Z[X]$，且以 $\alpha^{-1}$ 为根，故 $\alpha\in\mathcal O_K^\times$.
\end{proof}
\subsection{挠子群}\label{sec:nt-5-3-2}
由挠元素组成的\newterm{nt-torsion-subgroup}{挠子群}{torsion subgroup} $(\mathcal O_K^\times)_{\mathrm{tors}}$，就是 $K$ 中所有单位根组成的集合 $\mu(K)$.
\begin{lemma*}
$\mu(K)$ 是有限群，因而是循环群.
\end{lemma*}
\begin{proof}
有限性由分圆多项式在 $\mathbb Q$ 上的不可约性推出.事实上，$\zeta_n\in\mu(K)\Rightarrow\varphi(n)\mid[K:\mathbb Q]$.循环性则是一个经典结果：域的乘法群的每个有限子群都是循环群.
\end{proof}
\begin{remark*}
若 $K$ 有实嵌入（即 $r_1\ne0$），则 $\mu(K)=\{\pm1\}$.
\end{remark*}
\begin{example*}[二次域]
若 $\zeta_n\in\mu(\mathbb Q[\sqrt d])$，则 $\varphi(n)\le2$，所以 $n$ 的可能值为 $2,3,4$ 或 $6$.\footnote{译注：原文列举中略去了平凡单位根对应的 $n=1$；这不影响下面的单位根群分类.}此外，容易看出 $\zeta_3\in\mathbb Q[\sqrt d]\Longleftrightarrow d=-3$，且 $\zeta_4\in\mathbb Q[\sqrt d]\Longleftrightarrow d=-1$.由此得到
\begin{align*}
\mu(\mathbb Q[\sqrt d])=
\begin{cases}
\mu_6&\text{若 }d=-3,\\
\mu_4&\text{若 }d=-1,\\
\mu_2&\text{若 }d\ne-1,-3.
\end{cases}
\end{align*}
\end{example*}
\begin{proposition*}\label{prop:nt-5-3-2-modulus-one}
设 $\alpha\in K$，且对每个 $\sigma\colon K\hookrightarrow\mathbb C$ 都有 $|\sigma(\alpha)|=1$.则 $\alpha\in\mu(K)$.\footnote{译注：原命题缺少 $\alpha$ 为代数整数的假设.例如 $\alpha=(3+4i)/5\in\mathbb Q(i)$ 的所有共轭的模均为 $1$，但它不是单位根.下面证明中将 $\alpha$ 的各次幂放入由代数整数组成的 $E_{[K:\mathbb Q],1}$，需要这一假设.这里保留原陈述.}
\end{proposition*}
\begin{proof}
更一般地，有如下结果：设 $m,M>0$，则下面的集合是有限的：
\begin{align*}
E_{m,M}:=\{\alpha\in\overline{\mathbb Z},\ \deg(f_\alpha)\le m\text{ 且 }|\sigma(\alpha)|\le M,\ \forall\sigma\colon K\hookrightarrow\mathbb C\}.
\end{align*}
这是显然的，因为这时 $f_\alpha$ 的整数系数均可由 $M$ 控制而有界.现在只须注意，若 $\alpha$ 如命题所述，则 $\{\alpha^n,\ n\in\mathbb N\}\subset E_{[K:\mathbb Q],1}$.
\end{proof}
\subsection{自由部分}\label{sec:nt-5-3-3}
Abel 群 $\mathcal O_K^\times$ 的秩，按定义是 $\mathbb Q$ 向量空间 $\mathbb Q\otimes_{\mathbb Z}\mathcal O_K^\times$ 的维数.事先并不清楚它是否有限.注意，这也是商群 $\mathcal O_K^\times/(\mathcal O_K^\times)_{\mathrm{tors}}$ 的秩，而此商群是无挠 Abel 群.
令 $\Sigma$ 如上一节所定义.定义映射
\begin{align*}
\mathcal L\colon K_{\mathbb R}^\times=\prod_{\sigma\in\Sigma}K_\sigma^\times&\longrightarrow\mathbb R^\Sigma,\\
(x_\sigma)_{\sigma\in\Sigma}&\longmapsto\bigl(\log(|x_\sigma|^{[K_\sigma:\mathbb R]})\bigr)_{\sigma\in\Sigma}.
\end{align*}
$\mathcal O_K^\times$（对角嵌入 $K_{\mathbb R}^\times$ 中）的像 $\mathcal L(\mathcal O_K^\times)$ 是 $\mathbb R^\Sigma$ 的加法子群，乘积公式说明它包含于超平面
\begin{align*}
H:=\left\{(a_\sigma)_{\sigma\in\Sigma}\in\mathbb R^\Sigma,\ \sum_{\sigma\in\Sigma}a_\sigma=0\right\}.
\end{align*}
\begin{theorem*}[Dirichlet]\label{thm:nt-dirichlet-units}
沿用上述记号，有：\footnote{译注：原文 i) 的最后一项写为 $(\mathcal O_K)_{\mathrm{tors}}$，漏了乘法群的 $\times$；这里保留原式，其意为 $(\mathcal O_K^\times)_{\mathrm{tors}}$.}
\begin{enumerate}[label=\roman*)]
\item $\operatorname{Ker}(\mathcal L)\cap\mathcal O_K^\times=\mu(K)=(\mathcal O_K)_{\mathrm{tors}}$.
\item $\mathcal L(\mathcal O_K^\times)$ 是 $H$ 中的格（即离散且余紧）.
\end{enumerate}
特别地，$\mathcal O_K^\times$ 的秩有限，等于 $r_1+r_2-1$.
\end{theorem*}
\begin{proof}
\textup{i)} $\alpha\in\operatorname{Ker}(\mathcal L)\cap\mathcal O_K^\times\Longleftrightarrow|\sigma(\alpha)|=1$，$\forall\sigma$.根据\hyperref[prop:nt-5-3-2-modulus-one]{上面的命题}，这也等价于 $\alpha\in\mu(K)$.
\textup{ii)}
\begin{enumerate}[label=(\alph*)]
\item 先证明 $\mathcal L(\mathcal O_K^\times)$ 离散.若 $B$ 是 $H$ 中的有界开集，则 $\mathcal L^{-1}(B)$ 是 $K_{\mathbb R}^\times$ 中的有界开集，所以 $\mathcal L^{-1}(B)\cap\mathcal O_K$ 是有限集，因为 $\mathcal O_K$ 在 $K_{\mathbb R}^\times$ 中离散.于是 $B\cap\mathcal L(\mathcal O_K^\times)$ 有限.至此已知 $\mathcal O_K^\times$ 的秩有限.
\item 还需证明 $\mathcal L(\mathcal O_K^\times)$ 在 $H$ 中余紧.这将给出它的秩为 $\dim(H)=r_1+r_2-1$.只须证明 $\mathcal O_K^\times$ 在下述群中余紧：
\begin{align*}
\mathcal H:=\mathcal L^{-1}(H)
=\left\{x=(x_\sigma)_{\sigma\in\Sigma}\in K_{\mathbb R}^\times,\ N(x):=\prod_{\sigma\in\Sigma}|x_\sigma|^{[K_\sigma:\mathbb R]}=1\right\},
\end{align*}
这里的意思是：存在 $K_{\mathbb R}^\times$ 的紧集 $C$，使得 $\mathcal H=\mathcal O_K^\times(\mathcal H\cap C)$.
与定理~\ref{sec:nt-5-2-2} 的证明一样，对 $c=(c_\sigma)_{\sigma\in\Sigma}\in(\mathbb R_+^\times)^\Sigma$，考虑球的乘积
\begin{align*}
B(c):=\{x=(x_\sigma)_{\sigma\in\Sigma}\in K_{\mathbb R},\ |x_\sigma|\le c_\sigma,\ \forall\sigma\in\Sigma\},
\end{align*}
并选取 $c$ 使得 $\operatorname{vol}(B(c))=2^n\operatorname{covol}(\mathcal O_K)$.这就固定了 $N(c):=\prod_\sigma c_\sigma^{[K_\sigma:\mathbb R]}$ 的值（这个值具体是多少并不重要）.
现在，对 $x\in\mathcal H$，$x^{-1}B(c)$ 的体积为
\begin{align*}
N(x^{-1})\operatorname{vol}(B(c))=\operatorname{vol}(B(c))=2^n\operatorname{covol}(\mathcal O_K),
\end{align*}
所以 Minkowski 定理给出非零元素 $\alpha\in\mathcal O_K\cap x^{-1}B(c)$.于是有
\begin{align*}
|N_{K/\mathbb Q}(\alpha)|\le N(x^{-1})N(c)=N(c).
\end{align*}
因为 $\{I\subset\mathcal O_K,\ (\mathcal O:I)\le N(c)\}$ 是有限集，所以存在 $\alpha_1,\ldots,\alpha_N$ 使得\footnote{译注：下式原文量词使用 $a$，结论却写 $(\alpha)$；应理解为对同一元素所生成的主理想的陈述.原文也未在此量词中排除 $a=0$.这里保留原式.}
\begin{align*}
\forall a\in\mathcal O_K,\quad |N_{K/\mathbb Q}(a)|\le N(c)\ \Longrightarrow\ \exists j,\quad(\alpha)=(\alpha_j).
\end{align*}
因此存在 $j$ 使得 $(\alpha)=(\alpha_j)$，即 $\alpha\in\mathcal O_K^\times\alpha_j$，由此得到 $x\in\mathcal O_K^\times\alpha_j^{-1}B(c)$.于是令\footnote{译注：下式定义的 $C$ 在 $K_{\mathbb R}$ 中紧，但含有 $0$，故不是前面要求的 $K_{\mathbb R}^\times$ 中的紧集.可用 $C\cap\mathcal H$：它是 $K_{\mathbb R}$ 中的紧集且各分量均非零，从而也是 $K_{\mathbb R}^\times$ 中的紧集.这里保留原构造.}
\begin{align*}
C:=\bigcup_{j=1,\ldots,N}\alpha_j^{-1}B(c),
\end{align*}
这显然是紧集.我们已经证明 $\mathcal H=\mathcal O_K^\times\cdot(\mathcal H\cap C)$.\qedhere
\end{enumerate}
\end{proof}
定理推出 $\mathcal O_K^\times\simeq(\mathcal O_K^\times)_{\mathrm{tors}}\times\mathbb Z^{r_1+r_2-1}$.一个通常很困难的问题是求一组\newterm{nt-fundamental-units}{基本单位}{fundamental units}，即一族单位 $(\varepsilon_1,\ldots,\varepsilon_{r_1+r_2-1})$，使它们在 $\mathcal O_K^\times/(\mathcal O_K^\times)_{\mathrm{tors}}$ 中的像构成一组 $\mathbb Z$-基.
\begin{example*}
在 $\mathbb Q(\sqrt d)$、$d>0$ 的情形，自由部分的秩为 $1$.当 $\mathcal O_{\mathbb Q(\sqrt d)}=\mathbb Z[\sqrt d]$ 时，单位 $\varepsilon=x+y\sqrt d$ 满足 Pell--Fermat 方程 $N(\varepsilon)=x^2-dy^2=\pm1$，该方程的求解用到 $\sqrt d$ 的连分数展开.
\end{example*}
\chapter*{习题}
\addcontentsline{toc}{chapter}{习题}
\label{chap:nt-exercises}
\begin{exercise*}[证明范数的传递性]\label{ex:nt-norm-transitivity}
给定环 $B$ 上有限秩自由模 $M$ 的自同态 $u$，其中 $B$ 本身是 $A$ 上的有限秩自由模.要证明 $A$ 中的等式
\begin{align*}
\mathcal E(A,B,M,u)\colon\qquad\det_A(u)=N_{B/A}(\det_B(u)).
\end{align*}
\begin{enumerate}[label=\roman*)]
\item\label{item:nt-norm-transitivity-i} 当 $B=A\times A\times\cdots\times A$ 时，验证这个等式.
\item\label{item:nt-norm-transitivity-ii} 设 $f\colon A\longrightarrow A'$ 是环同态.令 $B':=A'\otimes_A B$、$M':=A'\otimes_A M=B'\otimes_B M$，并令 $u':=\operatorname{id}\otimes u$.证明：
\begin{enumerate}[label=(\alph*)]
\item\label{item:nt-norm-transitivity-ii-a} $\mathcal E(A,B,M,u)\Rightarrow\mathcal E(A',B',M',u')$.
\item\label{item:nt-norm-transitivity-ii-b} 若 $f$ 为单射，则 $\mathcal E(A',B',M',u')\Rightarrow\mathcal E(A,B,M,u)$.
\end{enumerate}
\item\label{item:nt-norm-transitivity-iii} 利用\hyperref[item:nt-norm-transitivity-i]{i)} 和\hyperref[item:nt-norm-transitivity-ii-b]{ii)(b)}，证明当 $A$ 是域、$B$ 是 $A$ 上的可分代数时，$\mathcal E(A,B,M,u)$ 成立.再由此推出 $B$ 为特征零整环的情形（这基本上就是本课程所需要的全部情形）.
\item\label{item:nt-norm-transitivity-iv} 现在处理 $B=A[X]/(f)$、$f\in A[X]$ 首一的情形.
\begin{enumerate}[label=(\alph*)]
\item\label{item:nt-norm-transitivity-iv-a} 证明下面的特殊情形可用\hyperref[item:nt-norm-transitivity-iii]{iii)} 处理：
\begin{align*}
A&=\mathbb Z[\alpha_1,\ldots,\alpha_m]\bigl[a_{ij}^{(0)},\ldots,a_{ij}^{(m-1)}\bigr]_{1\le i,j\le n},\\
f&=X^m+\alpha_1X^{m-1}+\cdots+\alpha_m.
\end{align*}
\item 利用\hyperref[item:nt-norm-transitivity-ii-a]{ii)(a)}，通过特殊化论证，由此推出\hyperref[item:nt-norm-transitivity-iv]{iv)} 的一般情形.在\hyperref[item:nt-norm-transitivity-iv-a]{iv)(a)} 的情形中，可以取 $M=B^n$，并令 $u$ 由矩阵 $(b_{ij})_{i,j}$ 给出，其中 $b_{ij}=\sum_{k=0}^{m-1}a_{ij}^{(k)}\overline X^k$.
\end{enumerate}
\item 由\hyperref[item:nt-norm-transitivity-iv]{iv)} 推出，当 $B$ 是域时 $\mathcal E(A,B,M,u)$ 成立（可以考虑 $A$ 与 $B$ 之间的单生成中间扩张）.再由此推出 $B$ 为整环的情形.
\end{enumerate}
\end{exercise*}
